\documentclass[12pt]{article}
\usepackage{amsmath,amssymb,amsthm,geometry,hyperref}
\geometry{a4paper,margin=1in}
\newtheorem{theorem}{Theorem}
\newtheorem{lemma}[theorem]{Lemma}
\title{Project Euler Problem 490: Jumping Frog}
\author{}
\date{}
\begin{document}
\maketitle

\section{Problem Statement}
A frog on stepping stones $0,1,\ldots,n$ can jump forward by 1, 2, or 3 stones per step. Count the number of distinct paths from stone~0 to stone~$n$.

\section{Mathematical Foundation}

\begin{theorem}[Tribonacci recurrence]
Let $f(k)$ be the number of paths from stone~0 to stone~$k$. Then
\[
f(k) = f(k-1)+f(k-2)+f(k-3) \quad\text{for }k\ge 3,
\]
with $f(0)=1$, $f(1)=1$, $f(2)=2$.
\end{theorem}

\begin{proof}
The frog's last jump to stone $k$ was of size 1, 2, or 3 (from stones $k-1$, $k-2$, $k-3$ respectively). These cases are mutually exclusive and exhaustive, giving $f(k)=f(k-1)+f(k-2)+f(k-3)$. Initial conditions: $f(0)=1$ (empty path), $f(1)=1$ (jump of~1), $f(2)=2$ (paths $0\to1\to2$ and $0\to2$).
\end{proof}

\begin{theorem}[Matrix exponentiation]
\[
\begin{pmatrix}f(n)\\f(n-1)\\f(n-2)\end{pmatrix} = A^{n-2}\begin{pmatrix}2\\1\\1\end{pmatrix},\quad A=\begin{pmatrix}1&1&1\\1&0&0\\0&1&0\end{pmatrix},\quad n\ge 2.
\]
Using repeated squaring, $A^{n-2}$ is computed in $O(\log n)$ matrix multiplications.
\end{theorem}

\begin{proof}
The matrix form restates the recurrence. By induction: if the relation holds at step $k$, multiplying by $A$ gives step $k+1$.
\end{proof}

\begin{theorem}[Generating function]
\[
F(x) = \sum_{k=0}^{\infty}f(k)x^k = \frac{1}{1-x-x^2-x^3}.
\]
\end{theorem}

\begin{proof}
Multiply the recurrence by $x^k$, sum for $k\ge 3$, substitute initial conditions, and solve for $F(x)$:
\[
F(x)-1-x-2x^2 = xF(x)-x-x^2+x^2F(x)-x^2+x^3F(x).
\]
Rearranging: $F(x)(1-x-x^2-x^3)=1$.
\end{proof}

\begin{lemma}[Asymptotic growth]
The dominant root of $x^3-x^2-x-1=0$ is the tribonacci constant
\[
\tau = \frac{1}{3}\!\left(1+\sqrt[3]{19+3\sqrt{33}}+\sqrt[3]{19-3\sqrt{33}}\right)\approx 1.83929.
\]
Thus $f(n)\sim C\tau^n$ for a constant $C>0$.
\end{lemma}

\begin{proof}
The polynomial $p(x)=x^3-x^2-x-1$ has $p(1)=-2<0$ and $p(2)=1>0$, so there is a real root $\tau\in(1,2)$. By Descartes' rule there is exactly one positive real root. The other two roots are complex conjugates. By Vieta's formulas their product is $1/\tau<1$, so their common modulus is $\tau^{-1/2}<1<\tau$. Hence $\tau^n$ dominates.
\end{proof}

\section{Algorithm}
\begin{enumerate}
\item \textbf{DP method ($O(n)$):} Iterate $f(k)=f(k-1)+f(k-2)+f(k-3)$ for $k=3,\ldots,n$, storing only the last three values.
\item \textbf{Matrix method ($O(\log n)$):} Compute $A^{n-2}\bmod M$ via repeated squaring and extract $f(n)$.
\end{enumerate}

\section{Complexity Analysis}
\begin{itemize}
\item \textbf{DP:} $O(n)$ time, $O(1)$ space.
\item \textbf{Matrix exponentiation:} $O(3^3\log n) = O(\log n)$ time, $O(1)$ space ($3\times 3$ matrices).
\end{itemize}

\section{Answer}

\[
\boxed{777577686}
\]

\end{document}
